\chapter{Conclusiones}
\label{cap:conclusiones}

\section{Discusión}

Los módulos anteriores comparten una lección de método. El diseño inicial trataba la \emph{estimación de propensión} y la \emph{recomendación} como un mismo problema y asumía, sin contrastarlo, que el perfil demográfico del usuario sería la variable determinante. Separar ambos problemas y exigir que fuese la evidencia, y no la intuición de partida, quien decidiese qué señal importa en cada uno cambió el diagnóstico. Qué señal resulta finalmente útil para cada tarea se detalla en la sección siguiente; lo relevante aquí es que esa pregunta se respondió midiendo, no suponiendo.

La segunda lección es que la forma de validar puede pesar más que la elección del modelo. La fuga de datos en el módulo de propensión no nacía de un fallo del algoritmo, sino de un esquema de validación que dejaba al mismo usuario en entrenamiento y prueba; corregirla rebajó el rendimiento aparente del modelo demográfico a su valor realista. En el recomendador ocurrió algo parecido: reconstruir el segmento usando solo datos de entrenamiento eliminó una ventaja que era un artefacto. Comparar cada modelo contra \textit{baselines} triviales y respaldar cada afirmación con tests pareados (bootstrap, McNemar, Wilcoxon) y corrección por comparaciones múltiples evitó dar por buenas mejoras inexistentes.

La tercera lección es que, en este problema, lo que más limita el rendimiento no es el modelado sino los datos. Las variables demográficas del usuario son sintéticas, muestreadas de distribuciones agregadas, y aportan poca señal individual; ese techo acota lo que cualquier algoritmo puede lograr y condiciona la lectura de todos los módulos.

\begin{hallazgo}
\textbf{Conclusión central.} La señal útil para la \emph{propensión} reside en el contenido del producto
(sector, categoría y texto del asunto); la útil para \emph{recomendar} reside en el \emph{comportamiento}
del usuario, no en su perfil demográfico. Las variables demográficas sintéticas marcan el techo del
rendimiento, por lo que el valor del trabajo está en la metodología.
\end{hallazgo}

\section{Conclusiones principales}

Este trabajo ha desarrollado un \textit{framework} integral de aprendizaje automático para la optimización de campañas de email marketing en el mercado español, cubriendo el ciclo completo desde la ingesta de datos brutos hasta la recomendación personalizada, la previsión de demanda y un ejemplo de despliegue. Más allá de los módulos, la principal aportación es metodológica: un pipeline reproducible que toma cada decisión de forma justificada (validación sin fuga, comparación frente a \textit{baselines} triviales antes de adoptar un modelo, corrección del sesgo de muestreo y contraste estadístico) y que distingue con evidencia qué señal es útil para cada problema, qué modelo merece desplegarse y a qué coste. Su valor no reside en una métrica de exactitud máxima, sino en un proceso disciplinado que reporta también lo que \emph{no} funciona. Leído así, el trabajo funciona como una guía para desarrollar un proyecto de datos orientado a maximizar el ROI: cada etapa muestra qué construir, cómo validarlo y cómo fundamentar la decisión de desplegar o no, incluida la de \emph{no} lanzar una campaña cuando la economía del canal no lo justifica.

De las cuatro hipótesis falsables planteadas en la introducción, se \emph{confirma} \textbf{H1} (las variables de producto mejoran significativamente la propensión, $p\approx0$) y se \emph{rechazan} las otras tres: H2 (el estudio de ablación muestra que las demográficas sintéticas no aportan poder predictivo), H3 (el recomendador por cluster no supera a la popularidad) y H4 (Prophet no mejora a los baselines ingenuos). Tres hipótesis rechazadas no son un fracaso: delimitan con evidencia dónde reside, y dónde no, la señal útil del problema.

Propensión y recomendación son problemas distintos, y la señal útil difiere según el objetivo. Para estimar la \emph{propensión} al clic, las variables de \emph{producto} (el sector, la categoría y, sobre todo, el texto del asunto del email codificado mediante \textit{embeddings}) aportan de forma estadísticamente significativa: el modelo con variables de producto (PR-AUC 0,402) supera tanto al modelo basado solo en variables de usuario como a los baselines de popularidad y de tasa de clic por categoría (véase el Capítulo~\ref{cap:propension}), lo que confirma que no se limita a memorizar la propensión media de cada categoría, sino que explota el texto del asunto. Una comparación de cinco algoritmos y dos arquitecturas (plano frente a jerárquico) confirma que el modelo plano supera al jerárquico en todos los algoritmos y que la red neuronal no aporta en datos tabulares con señal débil; el techo de rendimiento ($\approx0{,}40$ de PR-AUC) lo marcan los datos, no el algoritmo. Para \emph{recomendar}, en cambio, lo que discrimina es el \emph{comportamiento} del usuario (su segmento), no su perfil demográfico.

La validación cambia las conclusiones. La detección y corrección de una fuga de datos en M2 (el uso de validación cruzada por evento en lugar de agrupada por usuario) reveló que el rendimiento real del modelo demográfico (AUC-ROC 0,563) era muy inferior al inicialmente estimado (0,695). Este hallazgo constituye una contribución metodológica: ilustra cómo una elección de validación inadecuada puede inflar de forma engañosa las métricas en problemas donde las features son atributos estables del usuario.

El recomendador por cluster aporta valor operativo, no predictivo. Sirve sus recomendaciones \emph{a nivel de cluster} (segmento comportamental M1). Al evaluarlo sin fuga, reconstruyendo el segmento solo con datos de entrenamiento, \emph{no supera} a la popularidad en ninguna métrica (HR@5 0,586 frente a 0,629 a favor de la popularidad), un resultado confirmado por un test de Friedman seguido de Wilcoxon con corrección de Holm; la ventaja aparente del segmento contaminado era un artefacto de fuga temporal. Algo análogo ocurre con el cluster demográfico de M0: es la variable más \emph{usada} por M2 (3.ª por número de \textit{splits}), pero el estudio de ablación demuestra que su aporte predictivo es nulo, ya que prescindir de él y del resto de variables de usuario no reduce el PR-AUC. El valor de ambos es, por tanto, \emph{operativo}: segmentos interpretables, \textit{ranking} homogéneo por grupo y solución natural de arranque en frío, asignando a cada usuario nuevo las afinidades de su grupo. No mejoran el \textit{ranking} sobre una popularidad que, con un catálogo de 25 categorías y clics concentrados, resulta un \textit{baseline} muy fuerte.

El módulo M4 produce predicciones en el rango 259--452 clics/semana para enero--febrero 2026, con intervalos de confianza amplios coherentes con un histórico activo de solo 22 semanas. Una re-auditoría mediante \textit{backtest} rolling de 1 paso revela, además, que Prophet no supera a baselines ingenuos (naïve y media móvil): su error es mayor (MAE 750 frente a 502) aunque, con solo 14 orígenes y la corrección de muestra pequeña de Harvey-Leybourne-Newbold, la diferencia no alcanza significación al 5\% (Diebold-Mariano $p\approx0{,}06$), por lo que la evidencia es sugestiva, no concluyente. En una serie corta y volátil, ``la próxima semana será como la última'' iguala o supera a Prophet, un resultado negativo que cuestiona su valor añadido en este contexto.

Para usuarios con historial, un modelo \emph{warm} que añade su comportamiento pasado mejora la propensión de forma notable y consistente (PR-AUC $0{,}83\,\to\,0{,}94$, con mejora en los cinco cortes de una validación cruzada temporal), frente al modelo demográfico \emph{cold} válido para usuarios nuevos. Las probabilidades del modelo plano están razonablemente calibradas (\textit{Brier} 0,167 en la escala de entrenamiento). El modelo plano de propensión se expone en la API de ejemplo mediante un \textit{endpoint} \texttt{/propensity} que, dado un usuario conocido o un perfil nuevo, devuelve la probabilidad de clic desglosada por sector y por producto (corregida al prior real).

El análisis económico, interpretando el CPL como beneficio por clic, muestra que, bajo coste marginal por email, priorizar los envíos con el modelo mejora el ROI en las tres estructuras de coste analizadas (fijo, variable y mixto; por ejemplo, de 101\,\% a 148\,\% en el caso fijo), descartando los envíos cuyo beneficio esperado no cubre su coste y rescatando a positivo incluso los escenarios donde enviar a todos es deficitario. El punto de operación óptimo depende del objetivo de negocio (eficiencia frente a alcance). Ahora bien, esa comparación se hace sobre la muestra balanceada; al recalcular el ROI a escala de producción (reponderando al prior real del 2\,\% y decidiendo con la probabilidad corregida), el panorama es más crudo: enviar a todos pierde dinero en todas las estructuras y ni siquiera el mejor caso (coste fijo bajo) llega a ser rentable, pues el modelo lo lleva al punto de equilibrio ($\approx-1$\,\%, frente al $-$83\,\% de enviar a todos) con una fracción mínima de envíos. La lectura de negocio es que, a tasas de clic realistas, el targeting es imprescindible pero la rentabilidad la dicta la economía unitaria (coste por email frente a CPL).

\section{Limitaciones}

\begin{enumerate}
    \item \textbf{Variables sintéticas.} Las variables demográficas enriquecidas (situación laboral, estado civil, tamaño del hogar, habitaciones y tenencia de vehículo) son imputaciones estocásticas a partir de distribuciones agregadas del INE; no se observan a nivel individual. Un estudio de ablación lo confirma cuantitativamente: un modelo de propensión con \emph{solo} variables de producto iguala al modelo completo, mientras que uno con \emph{solo} variables de usuario apenas supera el azar (PR-AUC 0,278 vs.\ 0,238). Esto fija un techo bajo a la señal demográfica y explica por qué tanto el modelo inverso como el recomendador demográfico tienen poca capacidad discriminativa. Es la limitación dominante del trabajo.

    \item \textbf{Esparsidad extrema.} El 83,8\% de los usuarios tiene un único evento, lo que impide construir perfiles comportamentales fiables para la mayoría de la base y explica que el componente de historial individual del recomendador no aporte (peso $\alpha=0$).

    \item \textbf{Catálogo reducido.} Con solo 25 categorías de producto y clics concentrados en unas pocas, el baseline de popularidad es muy fuerte, lo que limita el margen de mejora medible de cualquier recomendador.

    \item \textbf{Tamaño de la subpoblación \textit{warm}.} El modelo \emph{warm} solo aporta a los usuarios con historial, que son minoría; para la mayoría (un evento) sigue rigiendo el límite de la señal demográfica sintética.

    \item \textbf{Previsión con histórico corto.} Prophet dispone de solo 22 semanas de actividad, lo que limita la fiabilidad del forecast más allá de pocas semanas y produce intervalos de confianza amplios. M4 se re-auditó con un \textit{backtest} rolling y el test de Diebold-Mariano (con corrección de muestra pequeña de Harvey-Leybourne-Newbold), pero la escasez de orígenes (14) limita la potencia estadística, por lo que la evidencia es sugestiva.

    \item \textbf{El modelo capta solo una etapa del embudo (clic dado apertura).} Un correo atraviesa dos filtros sucesivos: envío\,$\rightarrow$\,apertura\,$\rightarrow$\,clic. Como los datos solo contienen correos \emph{abiertos}, la variable objetivo es ``clic \emph{dado que} se abrió'': el modelo estima $P(\text{clic}\mid\text{apertura})$ y no modela la primera etapa (la probabilidad de que el correo llegue a abrirse). En consecuencia, la corrección de \textit{prior} al 2\,\% y el ROI ``por email enviado'' son exactos solo si la tasa de apertura es parecida entre productos; si un producto se abre mucho más que otro, su propensión real por envío diferiría aunque su clic-dado-apertura fuese idéntico. La extensión natural es un modelo en dos fases (uno para la apertura y otro para el clic) cuyo producto, $P(\text{abre})\cdot P(\text{clic}\mid\text{abre})$, recupera la propensión real por email enviado.
\end{enumerate}

\section{Trabajo futuro}

\begin{enumerate}
    \item \textbf{Datos reales en lugar de imputados.} La mejora con mayor impacto potencial es sustituir las variables demográficas sintéticas por observaciones reales (CRM, comportamiento en la web del anunciante). Todo el pipeline rendiría más sin cambios en el modelado.

    \item \textbf{Profundizar el modelo \textit{warm}.} El modelo \emph{warm} básico ya se implementó y evaluó (mejora demostrada para usuarios con historial). Como extensión, incorporar más señales de comportamiento (secuencia temporal de interacciones, recencia por sector) y un esquema de actualización incremental.

    \item \textbf{Mejorar la previsión temporal.} La re-auditoría mostró que Prophet no bate a baselines ingenuos en esta serie. Como alternativa, explorar modelos específicos para series cortas (suavizado exponencial, ARIMA con diagnóstico) o incorporar covariables de campaña (nº de envíos, presupuesto) como regresores externos.

    \item \textbf{Evaluación en producción e incrementalidad.} La validación \textit{offline} tiene límites conocidos en recomendación; un test A/B controlado (con grupo de envío y grupo de control) mediría el impacto real sobre la tasa de clic y el retorno de inversión, y permitiría estimar la \emph{incrementalidad} (\textit{uplift}) de cada envío, es decir, su efecto causal, que el análisis de ROI no captura al medir propensión y no efecto causal.
\end{enumerate}
